% musr_report.tex — MuSR 实验报告(与 musr_report.md 逐节对应)
% 编译:xelatex musr_report.tex(跑两遍);字体:Droid Sans Fallback(系统自带)
% 版式沿用 belief_r_report.tex(同一 preamble):md 的引用块/代码块在 PDF 里落为底色框,
% 正文文字对齐 md,改动全在呈现层(配色/框/表头/页眉/标题)。
\documentclass[fontset=none,10.5pt]{ctexart}
\setCJKmainfont[AutoFakeBold=2.5]{Droid Sans Fallback}
\setCJKsansfont[AutoFakeBold=2.5]{Droid Sans Fallback}
\setCJKmonofont{Droid Sans Fallback}
% U+2014(——)不在 Droid Sans Fallback 内,改走西文字体(Latin Modern 有该字形)
\xeCJKDeclareCharClass{Default}{"2014}
\usepackage[a4paper,margin=1.85cm]{geometry}
\usepackage{graphicx,booktabs,amssymb,amsmath,caption,enumitem,array,float,titlesec}
\usepackage{xcolor,colortbl,fancyvrb,fancyhdr,lastpage,microtype}
\usepackage[most]{tcolorbox}

% ---- 配色(克制:一个主色 + 两档灰,不上第二种彩色)----
\definecolor{accent}{HTML}{1F4E79}       % 主色:深学术蓝
\definecolor{accentsoft}{HTML}{EDF2F8}   % 提示框底色
\definecolor{codebg}{HTML}{F5F5F2}       % 代码框底色(暖灰)
\definecolor{codeframe}{HTML}{DCDCD6}
\definecolor{tabhead}{HTML}{EDF2F8}      % 表头底色
\definecolor{rulesoft}{HTML}{AFC0D2}
\definecolor{okgreen}{HTML}{2E7D4F}

\usepackage[colorlinks=true,linkcolor=accent,urlcolor=accent]{hyperref}

% ---- 紧凑排版 ----
\setlength{\parskip}{0.2em}
\setlength{\parindent}{0pt}
\emergencystretch=1.5em
\titlespacing*{\section}{0pt}{1.25ex plus .2ex}{0.6ex}
\titlespacing*{\subsection}{0pt}{0.75ex}{0.3ex}
\setlist{topsep=0.2em,itemsep=0.12em,parsep=0.12em,partopsep=0pt,leftmargin=1.6em}
\captionsetup{font=small,labelfont={bf,color=accent},skip=3pt}
\setlength{\textfloatsep}{7pt plus 2pt minus 2pt}
\setlength{\intextsep}{7pt plus 2pt minus 2pt}
\setlength{\abovecaptionskip}{3pt}
\setlength{\belowcaptionskip}{0pt}
\renewcommand{\arraystretch}{0.95}

% ---- 标题:主色 + 细分隔线 ----
\titleformat{\section}
  {\normalfont\large\bfseries\color{accent}}{\thesection}{0.6em}{}
  [\vspace{-0.7ex}{\color{rulesoft}\titlerule[0.7pt]}]
\titleformat{\subsection}
  {\normalfont\normalsize\bfseries\color{accent!88!black}}{}{0pt}{}

% ---- 页眉 ----
\pagestyle{fancy}
\fancyhf{}
\fancyhead[L]{\small\color{black!45}MuSR 实验报告}
\fancyhead[R]{\small\color{black!45}\thepage\,/\,\pageref{LastPage}}
\renewcommand{\headrule}{{\color{rulesoft}\hrule height 0.5pt}}
\renewcommand{\footrulewidth}{0pt}
\setlength{\headheight}{14pt}

% ---- 表格:表头底色,规则线不留缝 ----
\setlength{\aboverulesep}{0pt}
\setlength{\belowrulesep}{0pt}
\setlength{\extrarowheight}{2.2pt}
\arrayrulecolor{black!62}

% ---- md 里的灰框,在这里落成真的框 ----
% 引用块 → 左侧主色竖条 + 淡底
\newtcolorbox{notice}{
  enhanced, breakable, sharp corners,
  colback=accentsoft, colframe=accentsoft, boxrule=0pt,
  borderline west={2.4pt}{0pt}{accent!72},
  left=9pt, right=9pt, top=5.5pt, bottom=5.5pt,
  before skip=5pt, after skip=5pt,
  fontupper=\small,
}
% 代码块 → 暖灰底 + 细边框
\newtcolorbox{codebox}{
  enhanced, breakable, sharp corners,
  colback=codebg, colframe=codeframe, boxrule=0.5pt,
  left=8pt, right=8pt, top=5pt, bottom=5pt,
  before skip=5pt, after skip=5pt,
}

\begin{document}
\thispagestyle{empty}

\begin{center}
{\color{accent}\rule{\linewidth}{1.2pt}}\\[0.6em]
{\LARGE\bfseries\color{accent} MuSR 部分实验报告(论文 \S 5 全景)}\\[0.5em]
{\small\color{black!55} 2026-07-21}\\[0.55em]
{\color{accent}\rule{\linewidth}{0.5pt}}
\end{center}
\vspace{0.4em}

\section*{摘要}

本报告给出 \textbf{SAVI} 在 belief-tracking 任务上的实验结果。核心主张是:在没有免费验证器的真实任务上,SAVI 能修复``采样—投票''范式修不动的失败,且增益来自解码方式本身,不依赖额外信息。实验在信念追踪基准 MuSR object\_placements(Qwen3-4B)上分三步完成。(1) 搭建测试台:四级认证漏斗从 256 题中认证出 \textbf{19 道投票修不动的题}(另有 3 道知识对照);流行率测量说明同型失败步(全票写错,point mass)在未筛选题目上并不罕见,测试集严苛但不是孤例。(2) 排除池内替代解释:learned-verifier 重排与 best-of-N 逐题相等,gold-facts 掩码给出的池内上界也只有 3/19。``在采样池里挑得更聪明''复现不了修复,缺的候选不在池里;SAVI 的 support 轴(枚举候选)补的正是这一块。(3) 验证 SAVI 并归因:SAVI 实例化为枚举可达信念状态、每步锚一句模型自写的事件句、似然打分、Viterbi 取全局最优路径,\textbf{零人工标注修复 9/19},知识对照 0/3。消融把增益定位到``每步一句具体、成句断言''的存在(4$\to$9;内容改对仅 +1);机制探针显示似然打分特异、频率打分非特异,说明 SAVI 起效走的是``认得出''(recognize)通道。失败本身是分叉处的自信写错:给足信息后仍有 80\% 写错,不是缺信息。这也解释了修复为什么要由解码侧完成。

\section{实验目的}

\begin{enumerate}
\item \textbf{存在性:说明 commitment failure(承诺失败)在真实基准上存在。}认证出一个 can't 集:失败属于``知道但用不对'',而非知识缺失。(对应认证漏斗,第三节)
\item \textbf{转化:检验全局解码(枚举信念状态 + 模型自身似然打分 + 精确 Viterbi)能否让模型在这批题上实现 can't $\to$ can。}(对应第四、五节主结果:9/19 vs 投票 0/19)
\item \textbf{归因与定位:把转化的功劳归到具体成分(支撑 / 打分 / 每步锚句),并定位失败机制(分叉处自信地错)。}(对应四个 arm 消融 + 分叉重放,第五节)
\end{enumerate}

\section{Benchmark 及测试用例介绍}

\textbf{MuSR}(Sprague et al. 2024)是一个多步软推理基准,用 GPT-4 生成长叙事、人工校验答案,有三个子任务。本实验用 \textbf{object\_placements} 子任务,它是真实的信念追踪(belief tracking)基准:长叙事里物体被移动多次,每次移动只有部分角色看见;问``角色 X 认为物体 Y 在哪 / 会去哪找'',多选题。

实际的例子:

\begin{notice}
叙事:Ann 把钥匙放进\textbf{抽屉}后出门;她不在时 Ben 挪到\textbf{盒子}。\\
问题:Ann 会去哪找钥匙?\quad 选项:抽屉 / 盒子\\
正确答案是``抽屉''(她没看见挪动)——但模型顺着最近提到的「盒子」写下去,就回不了头。
\end{notice}

选择该任务的三个结构性原因:

\begin{itemize}
\item \textbf{状态可枚举}:每一步维护一张信念状态表——(观察者, \{物品: 位置\}),即每个观察者对每个物品位置的信念。表的取值有限,每步可达的候选状态能穷举出来。
\item \textbf{答案可直接读出}:终端信念表查 (X, Y) 一格就是答案,弃用模型自写的 ``ANSWER: \ldots'' 行,以排除``推到了正确终端状态、却把最后的答案行写错''这层混淆(读出消融单独验证过只 +1)。
\item \textbf{和 Countdown 的关键区别:域内没有免费的可执行检查器}。``免费''指同时满足三条:不需要训练、不需要真实事实/人工标注、从任务规则机械可算。Countdown 的检查器从域规则免费推出(执行算术求解),MuSR 的检查只能对照真实事实,三条至少违反一条。因此公平的强基线取多数投票。
\end{itemize}

{\sloppy
\textbf{主要实验模型与采样协议}:生成器 \textbf{Qwen3-4B};采样 T=1.0、top\_p=1.0、N 最大 1024、max\_new\_tokens=2048。\\
\textbf{主基线 = self-consistency(多数投票)}:N 条完整链各解析最后一行 \texttt{ANSWER: k},取答案众数。\par}

\section{实验设置:认证漏斗(如何选题)与基线定义}

后续主张都建立在``这 19 题确实修不动''上。总原则:每道门剔除一类并非真正修不动的题,对应如下:

\begin{center}\small
\begin{tabular}{@{}llp{11.5em}lp{13em}@{}}
\toprule
\rowcolor{tabhead}论文 & 内部阶段 & 操作 & 结果 & 排除的替代解释 \\
\midrule
起点 & — & object\_placements 256 题 & — & — \\
\textbf{F1} & 递增预算 & 错误众数题从 N=32 升到 N=1024 & \textbf{107 $\to$ 105},近乎持平 & ``预算不够'';预闭合后面所有比较的 token-parity 问题 \\
\textbf{F2} & 四门认证 & 对 105 道存活题过四道门 & \textbf{26 道认证} & 见下 \\
\textbf{F3} & 自提取控制 & 让模型自己提取事实再作答(两种注入方式) & 只修 5 道,\textbf{残余 21} & ``好 prompting 就够了'' \\
\textbf{F4} & 模板预检 & 结构化 prompt 下重跑投票 & 19/21 仍错,\textbf{最终 19} & ``结构化模板本身就能修'' \\
\bottomrule
\end{tabular}
\end{center}

\subsection{F1(budget escalation,256 $\to$ 105)}

\begin{itemize}
\item \textbf{判据}:多数投票的众数答案错(投票层面,不是单条链错),且预算从 N=32 加深到 N=1024 后众数仍错。结果 107@N=32 $\to$ 105@N=1024,近乎持平。
\item \textbf{隐含约束}:``不是预算不够''。投票加 32 倍 token 都不改答案,后续的比较不会再被``预算不足''这一解释推翻。
\end{itemize}

\subsection{F2(four checks,105 $\to$ 26)}

\begin{enumerate}
\item \textbf{统计稳定(排运气,剔 4)}:对``正确答案占比''与``错误众数占比''各计算 CI(N=1024),要求两区间分离(正确率上界 < 错误众数下界)才算真做不对。
\item \textbf{给定 gold-fact(排知识失败,剔 46)}:在 prompt 前注入观察日志(按时间逐条``谁把什么移到哪、谁看见了/没看见''),剔除注入后仍做不对的题目;这部分题目后续用作 knowledge control。
\item \textbf{不同规模过滤(排规模,剔 4)}:Qwen2.5-Instruct \{0.5, 1.5, 3, 7\}B;只要任一规模的模型``稳定解开''即剔除。
\item \textbf{鲁棒性检验(排种子巧合,剔 0)}:换第二种子,要求同一个错误选项仍为众数。
\end{enumerate}

\subsection{F3(self-extraction control,26 $\to$ 21)}

\begin{itemize}
\item \textbf{操作}:让模型先自己从叙事提取事实、再作答(两种注入方式)。只修好 5 道``易题''(其余题中位采样正确率 0.044)。
\item \textbf{记账语义}:F3/F4 修好的 5+2 道记为被剔除,不进任何方法成绩;能被好 prompting 修好的题不属于 can't 集。
\end{itemize}

\subsection{F4(template pre-check,21 $\to$ 19)}

\begin{itemize}
\item \textbf{操作}:在与解码侧共享的结构化模板(BELIEF[t] 信念表逐事件格式 + ANSWER 行)下重跑多数投票(N=64)。
\item \textbf{含义}:can't 不是格式伪影,投票基线与解码 prompt-matched:双方共享``信念表格式+任务框架'',解码独有的只剩单步条件化、事件句、机械读出。
\end{itemize}

\textbf{最终认证集构成}:19 题(+3 知识对照)= \textbf{15 道有覆盖}(采样到了、投票投错)+ \textbf{4 道零覆盖}(1024 次采样一次都没有对)。\\
\textbf{知识对照(knowledge controls)}:从 F2 (ii) 剔除的 46 道知识失败题(喂 gold-fact 仍错)里按种子抽样得到,用作``不该被修''的反向检验:任何方法把它们``修好'',就说明该修复对承诺失败不特异(可能泄露了答案),其修复数作废。

\subsection{流行率分析(prevalence,论文 \S 5.2)}

认证漏斗回答``哪些题修不动'',流行率回答``这种失败是不是罕见现象''。定义\textbf{全票写错步}(unanimous-wrong):该步所有采样链写出同一个错误后继,经验后继分布是 point mass(概率全压在一个点、零分歧)。结果:可判定题(共 40 道)中 \textbf{32/40 = 0.800} 至少含一个全票写错步(CI [0.675, 0.925]);可判定层 \textbf{67/115 = 0.583} 本身全票写错。\\
\textbf{与认证集选取的关系}:未筛选切片上,八成可判定题都带有与认证集同型的失败步(全票写错,主形态是照抄不更新),只是多数失败步没落在决定答案的位置。认证的 19 题就是失败恰好落在决定性分叉上、绕不过去的那部分题,不是筛选制造出来的罕见标本。

\section{实验组与对照组的设置}

\subsection{4.1 对照系:池内聚合(论文 \S 5.3)}

问题:同一批采样链,聚合得更聪明够不够?两个升级 arm 都基于结构化 prompt 的模型输出,且候选全部来自采样池,两个 arm 都不提供、不生成任何新候选状态:

\begin{itemize}
\item \textbf{learned-verifier}:训练一个步级事实一致性探针。用法:主 arm 是信念 trellis 上的 $\lambda$ 软加权解码;同一 verifier 另设 BoN 消融 arm(逐链取步级分均值、argmax 选链)。实测两者逐题完全相等(净差 0,CI[0,0]),即该 arm 等价于``逐步打分、选总分最高的链'';
\item \textbf{gold-facts mask}[ORACLE]:完全离线(post-hoc):先把 N 条链全部采完,解析每层 BELIEF 行 $\to$ 规范化 $\to$ 跨链归并成 trellis(归并率中位 0.963)$\to$ 用 gold-fact 一致性掩码禁掉矛盾结点 $\to$ 在剩余 trellis 上解码最优路径。机制上允许跨链拼接,实测拼接占比 0,修复全在单链内。
\end{itemize}

\textbf{本层结论}:learned-verifier = BoN(trellis 结构零增量);跨链拼接机制可用但没有被用上,修复只发生在完整采样链内;即使给 gold-fact 信号,池内上界也只有 3/19。问题不在挑选方式:缺的候选根本不在池子里。

\begin{center}\small
\begin{tabular}{@{}lll@{}}
\toprule
\rowcolor{tabhead}池内聚合方式 & fixed / 19 & 备注 \\
\midrule
majority voting & \textbf{0} & 主基线 \\
learned-verifier rerank(= best-of-N) & 净 +2,CI [0, 5] & 与 0 不可分;\textbf{逐题与 BoN 相等} \\
gold-facts mask (oracle), sampled chains only & 3 & 全在已有链内;跨链拼接 0 \\
\bottomrule
\end{tabular}

{\small\color{black!60}表 1\quad 对照系(池内聚合)结果}
\end{center}

\subsection{4.2 实验系:枚举解码 $\times$ 事件句四个 arm(论文 \S 5.4)}

\textbf{每步打分的模板}(确定性 teacher forcing,无采样;对每个 (层, 前状态, 候选) 三元组各算一次):

\begin{codebox}\small
\textbf{prefix}:\\
\textbf{[叙事全文]} + \textbf{[问题 + 选项]}\\
\textbf{[t-1 的信念状态]} BELIEF[t-1]: \{"Ben": \{"keys": "drawer"\}, ...\}\\
\textbf{[t event]}\; Event 3: the keys are moved to the box.\\
Write exactly one line recording what every character now believes ... in the same format.\\
+\\
\textbf{target}:\\
\textbf{[枚举可达状态]} BELIEF[t]: \{"Ben": \{"keys": "box"\}, ...\}\quad $\leftarrow$\; 打分只量这一段
\end{codebox}

\textbf{唯一被操纵的变量}:模板里的 \texttt{[event]} 事件句,\emph{``the keys are moved to the box.''},作为 anchor 锚定当前事件步:

\begin{center}\small
\begin{tabular}{@{}lp{24em}@{}}
\toprule
\rowcolor{tabhead}arm & 事件句内容 \\
\midrule
none & 整行删掉(只留一个空行) \\
placeholder & 固定占位句,语法完好、只透传步号 t,零事实内容 \\
\textbf{model-written} & 模型自己从叙事里提取(解析错了不剔除) \\
gold & 真实事实注释 \\
\bottomrule
\end{tabular}
\end{center}

\textbf{打分与归一化}:对 target 行逐 token log-softmax 求和得原始分;主归一化 = 同一 (层, 前状态) 候选集内的 log-softmax,即每个转移归一成候选集内的条件对数概率;路径分 = 逐层边权和,精确 Viterbi 取全局 argmax。\\
\textbf{对照}:除 3 道知识对照外,另设频率打分对照:同一枚举支撑上把似然换成采样频率打分,供 5.2 的 writes vs recognizes 区分使用。

\section{实验结果与分析}

\subsection{5.1 结果与分析}

\begin{center}\small
\begin{tabular}{@{}lll@{}}
\toprule
\rowcolor{tabhead} & fixed / 19 & + gold mask (ceiling) \\
\midrule
majority voting over the sampled pool & 0 & — \\
learned-verifier rerank of the pool(= best-of-N) & 净 +2,CI [0, 5] & 3 \\
\midrule
none & 3 & 5 \\
placeholder (no content) & 4 & 6 \\
\textbf{model-written (no human-labeled input)} & \textbf{9} & 12 \\
gold & 10 & 15 \\
\bottomrule
\end{tabular}

{\small\color{black!60}表 2\quad 主结果(各方法修复数 /19 与 gold-fact 掩码上界)}
\end{center}

\begin{itemize}
\item \textbf{机制}:具体断言让每步的后继分布坍缩;对全路径解码来说,``分布坍缩了''比``坍缩到哪个后继''更重要。
\item \textbf{这组梯度怎么读}:只有格式(占位句)几乎不起作用(none $\to$ placeholder 仅 3$\to$4);主要增益在``每步一句具体的、成句的断言''(placeholder $\to$ model-written,4$\to$9);断言内容改对只再值约一题(model-written $\to$ gold,9$\to$10)。
\item \textbf{自写句子大多有错,修复照样发生}:19 题没有一道全部提取正确,中位每题 3 个内容槽只对 1 个;event 内容质量与修复的相关弱(Spearman 0.31),有 6 道修复完全来自不精确的提取。
\item \textbf{仅枚举本身不起作用}:不依赖 event 句时,枚举几乎没有修复效果;gold arm 修复 4 道零覆盖题中的 3 道,删去 event 句后修复为 0。
\item \textbf{特异性}:似然打分下四个 arm 的知识对照都是 0/3。
\end{itemize}

\subsection{5.2 机制定位}

\begin{itemize}
\item \textbf{失败不是缺信息}:分叉重放中,给足正确历史和 gold-fact 句,认证集分叉处的众数后继只有 \textbf{20\%} 命中 gold-fact,知识对照为 \textbf{67\%}。模型是在分叉处自信地写错;认证阶段所说的``承诺失败''由此被定位到分叉节点。
\item \textbf{writes vs recognizes}:似然打分(recognize)特异(知识对照 0/3),频率打分(write)非特异(修 2/3 知识对照),起作用的是 recognize 一侧。另外,模型预测似然与实际采样频率的排序相关性 median $\rho$=0.400:TF 似然对模型实际采样频率的排序保真度不高,不能解读为模型的``信念''。
\end{itemize}

\section{思考与下一步计划}

\begin{enumerate}
\item \textbf{从信息角度分析增益(论文 \S 6 ``Same information, a different read-out'')}:SAVI 计算的每一个量都是采样器所读同一上下文的确定性函数。方法没有引入任何关于答案的新信息,增益只能落在提取泛函上:
  \begin{itemize}
  \item \textbf{support}:采样链只在它恰好写出的那一个后继上花费模型的判断;枚举把模型的条件判断读到每个可达信念状态上,包括它自己永远不会写出的反事实状态。仅有 trellis 结构没有帮助,但它是下游设计的必要条件;
  \item \textbf{scoring}:采样用模型写出的东西投票,而认证题恰恰是模型在分叉处自信写错的地方;似然打分改读``认得出'',并被每步锚句锐化(4$\to$9);
  \item \textbf{decision}:采样器每步承诺、不可逆,给出确定答案;解码器全局只承诺一次,过程中不锁死某个判断,答案从终端信念表机械读出(读出单独 +1)。
  \end{itemize}
  增益不来自信息量,而来自利用率:采样方法每条链的每一步都只把判断花在抽中的那一个后继上,SAVI 在每一步考虑全部可达后继,对同一信息的利用率更高。
\item \textbf{E0(aliasing audit 前置)}:E0 要求 BAR 定义在 beam baseline 的生成文本上。本实验最接近的量是 4.1 里 gold-facts mask 构建 trellis 时测得的跨链归并率中位 0.963:信号足够强(有大量内容可合并),但它来自 T=1.0 采样 + 结构化 BELIEF 模板,与 E0 的 beam 设定不能直接等价,需要扩展采样条件。
\item \textbf{BAR 预测搜索收益(Aliasing rate as a predictor of search benefit)}:在 4.1 的对照实验中,归并率高达 0.963,但 pool 内使用 SAVI 的 trellis 方式收益为零(与 BoN 逐题等价)。原因是 BAR 只测量``写出来的内容有多冗余'',观察不到``从没被写出来的状态'';决定性分叉上 pool 是 point mass,正确选项不在池里,归并得再好也选不出。因此要用 BAR 判断 SAVI 何时起作用,至少要补一个衡量决定性层覆盖 / 多样性的变量(point mass 检查,流行率的做法可迁移)。
\item \textbf{KV cache sharing}:目前实验中的归并率可以直接给 K'/K 提供 real-problem 数据点:0.963 意味着 state-level 实际束宽 K' 远小于 K,内存节省潜力大。这是纯粹的效率主张,与``trellis 对准确率零增量''的结论不冲突。
\item \textbf{上下文容量与推理能力}:模型的上下文容量和推理能力远不成比例:上下文能容纳很多步推理,但往往在步数还不大时出错频率就已经很高,且后续出错与带错的历史文本相关。与 MuSR 的关联:
  \begin{itemize}
  \item \textbf{符合前提}:认证题的失败发生在远离上下文上限的位置(叙事+链远小于 8192 token 上限,而 80\% 可判定题仍含全票错步)。出错不是``存不下'',是``用不好'',与``推理步数远小于容量时错误已高发''的观察一致;
  \item \textbf{状态化去噪}:走识别通道(枚举 + 似然),比较(规范化重置)与(不重置)两个对照的结果差异。
  \end{itemize}
\item \textbf{拓展实验}:在更多 benchmark(ZebraLogic 逻辑谜题、seqbench 迷宫问题等)上验证 SAVI 的效果。
\item \textbf{模型规模限制}:目前的模型规模($\le$7B)上可能不存在``belief''这一概念,计划验证 scale up 之后的效果。
\item \textbf{训一个可迁移的步级 verifier}:目前的 verifier 是任务特定的,在新任务上需要重新设计,迁移性差;考虑结合 LLM 入手。
\end{enumerate}

\end{document}
